\documentclass{article}
\usepackage{geometry}
\geometry{a4paper, margin=1cm}
\usepackage{iftex}
\ifXeTeX
\usepackage{fontspec}
\setmonofont{FreeMono}
\fi
\usepackage{minted}
\setminted{
    fontsize=\small,
    breaklines=true,
    linenos=false,
    autogobble=true,
    frame=single,
    framesep=10pt
}
\begin{document}
\tableofcontents

\section{Daniel's debug and development session}

\paragraph{Edit 1}
\textbf{File: \texttt{contacts.py}}

\begin{minted}{python}
"""
En modul som experimenterar med klasser.

Vi vill göra detta genom att skriva ihop en kontaktbok.
"""

import learnlog

class Contact:
    """
    Representerar en kontakt i kontaktboken.
    Håller koll på alla data som rör den personen.
    """
    def __init__(self, name, phone=None, email=None, address=None):
        """
        Tar alla attribut som kontakten ska ha.
        """
        self.name = name
        self.phone = phone # contact.phone = value
        self.email = email
        self.address = address

    # a = contact.phone
    @property
    def phone(self):
        ''' method that returns the phone number when contact.phone is used '''
        return self.__phone

    # contact.phone = '070187237'
    @phone.setter
    def phone(self, value):
        '''  sets phone number to a new value if it's a number'''

        # TODO: edit so value can be set to None
        if  value is not None or not value.isdigit():
            raise ValueError("Telefonnummret får bara innehålla siffror")
        self.__phone = value

    @property
    def email(self):
        ''' method that returns the email when contact.email is used'''
        return self.__email

    @email.setter
    def email(self, value):
        if '@' not in value:
            raise ValueError("E-postadressesn måste innehålla @")
        self.__email = value


    def __lt__(self, other):
        ''' Om self < other returnerar vi True annars False'''
        if self.name < other.name:
            return True
        elif self.name == other.name and self.email < other.email:
            return True
        elif self.name == other.name and self.email == other.email \
                and self.address < other.address:
            return True
        elif self.name == other.name and self.email == other.email \
                and self.address == other.address and self.phone < other.phone:
            return True
        return False

class ContactBook:
    ''' klass som innehåller kontaker, representerar en vanlig kontaktbok'''
    def __init__(self):
        ''' metod som skapar en lista med kontakter när ett kontaktboksobjekt skapas'''
        self.__contacts = []

    @property
    def contacts(self):
        ''' returnerar en kopia av kontaktlistan contacts när contactbook.contacts används'''
        return self.__contacts.copy()

    def add_contact(self, contact):
        ''' lägger till en kontakt till kontaktlistan om den ät av typ Contact'''
        if not isinstance(contact, Contact):
            raise TypeError("Kontakten måste vara av typen Contact")
        self.__contacts.append(contact)

    def find_contact(self, search_terms):
        ''' '''
        num_matches = {}
        for contact in self.__contacts:
            for search_term in search_terms.split():
                if match_contact(search_term, contact):
                    # räkna upp antal matchningar
                    if contact not in num_matches:
                        num_matches[contact] = 1
                    else:
                        num_matches[contact] += 1

        #print(num_matches)
        # sortera kontakterna i num_matches efter antal matchningar.
        # num_matches = {kontakten: antalet, kontakten: antalet, ...}
        # num_match_tuples = [(antalet, kontakten), (antalet, kontakten), ...]
        num_match_tuples = []
        for contact, num in num_matches.items():
            num_match_tuples.append(Match(num, contact))

        return num_match_tuples

def add_contact(contactbook):
    """
    Hanterar menydelen att lägga till en kontakt.
    contacts är en lista med kontakter.
    """
    name = input("Ange fullständigt namn för ny kontakt: ")
    contact = Contact(name)
    while phone := input("Ange telefonnummer för ny kontakt: "):
        try:
            contact.phone = phone
            break
        except ValueError as err:
            print(f"{err}: Försök igen!")

    while email := input("Ange email för ny kontakt: "):
        try:
            contact.email = email
            break
        except ValueError as err:
            print(f"{err}: Försök igen!")

    address = input("Ange adress för ny kontakt: ")
    contact.address = address

    contactbook.add_contact(contact)

def find_contact(contactbook):
    """
    Hanterar menydelen för att söka efter kontakter.
    """
    while True:
        search_terms = input("Ange sökterm: ").strip()

        num_match_tuples = contactbook.find_contact(search_terms)

        sorted_matches = sorted(num_match_tuples)
        # loopa igenom den sorterade listan och skriv ut varje kontakt.
        for match in sorted_matches:
            print_contact(match.contact)
            print()

class Match:
    ''' Ett sökresultat/ en matchning, där num är antalet matchade söktermer i 
    en kontakt och conatct är kontakten det gäller'''
    def __init__(self, num, contact):
        self.num = num
        self.contact = contact

    def __lt__(self, other):
        '''a och be är tuplar (num, contact), returnerar sant om num i a är större än num i b, 
        om båda är lika returnerar vi sant om contact i a är mindre än contact i b'''
        # self < other

        if self.num > other.num:
            return True
        elif self.num == other.num and self.contact < other.contact:
            return True
        return False

def match_contact(search_term, contact):
    """
    Returnerar True om något av attributen hos contact
    matchar söktermen (search_term).

    Med matchar menar vi att search_term är en delsträng.
    Telefonnummer använder dock exakt jämförelse.
    """
    if search_term.casefold() in contact.name.casefold():
        return True
    elif search_term.casefold() == contact.phone:
        return True
    elif search_term.casefold() in contact.email.casefold():
        return True
    elif search_term.casefold() in contact.address.casefold():
        return True

    return False


def print_contact(contact):
    """
    Skriver ut datat för en kontakt.
    """
    print(f"  Namn: {contact.name}")
    print(f"  Tele: {contact.phone}")
    print(f"  Mail: {contact.email}")
    print(f"  Adr : {contact.address}")

def main():
    """
    Huvudprogram
    """
    # vi vill kunna lägga till kontakter (kombinera med input)
    #               söka efter kontakter
    #               ändra kontakter
    #               skriva ut alla kontakter
    contactbook = ContactBook()
    contactbook.add_contact(Contact("Daniel", "0701233455", "dbosk@kth.se", "Någonannanstans"))
    contactbook.add_contact(Contact("Kawsar", "070", "kiaa@kth.se", "Någonstans"))

    while True:
        print()
        print("Välj ett alternativ (heltal):")
        print("1. Lägg till kontakt.")
        print("2. Sök efter kontakt.")
        print("3. Ändra en kontakt.")
        print("4. Skriv ut alla kontakter")
        print("5. Avsluta")
        choice = int(input("Mata in ditt alternativ: "))
        print()
        if choice == 1:
            add_contact(contactbook)
        elif choice == 2:
            find_contact(contactbook)
        elif choice == 3:
            pass
        elif choice == 4:
            pass
        elif choice == 5:
            return
        else:
            raise NotImplementedError("Inga andra alternativ är implementerade.")

if __name__ == "__main__":
    main()
\end{minted}

\noindent
\begin{minipage}[t]{0.48\linewidth}
\paragraph{Run 1}
Script: \texttt{contacts.py} \\
Result: Error \\
Exit code: 1

\begin{minted}{pytb}
ValueError: Telefonnummret får bara innehålla siffror
\end{minted}

\begin{minted}{text}
Traceback (most recent call last):
  File "/home/dbosk/teaching/prgi25/övn-2025-11-07/contacts.py", line 228, in <module>
    main()
    ~~~~^^
  File "/home/dbosk/teaching/prgi25/övn-2025-11-07/contacts.py", line 201, in main
    contactbook.add_contact(Contact("Daniel", "0701233455", "dbosk@kth.se", "Någonannanstans"))
                            ~~~~~~~^^^^^^^^^^^^^^^^^^^^^^^^^^^^^^^^^^^^^^^^^^^^^^^^^^^^^^^^^^^
  File "/home/dbosk/teaching/prgi25/övn-2025-11-07/contacts.py", line 19, in __init__
    self.phone = phone # contact.phone = value
    ^^^^^^^^^^
  File "/home/dbosk/teaching/prgi25/övn-2025-11-07/contacts.py", line 36, in phone
    raise ValueError("Telefonnummret får bara innehålla siffror")
ValueError: Telefonnummret får bara innehålla siffror
\end{minted}
\end{minipage}
\hfill
\begin{minipage}[t]{0.48\linewidth}
\paragraph{Edit 2}
\begin{minted}{diff}
--- a/contacts.py
+++ b/contacts.py
@@ -32,7 +32,7 @@
         '''  sets phone number to a new value if it's a number'''
 
         # TODO: edit so value can be set to None
-        if  value is not None or not value.isdigit():
+        if value is not None and not value.isdigit():
             raise ValueError("Telefonnummret får bara innehålla siffror")
         self.__phone = value
\end{minted}
\end{minipage}

\paragraph*{Variation-theory note}
\textit{Critical aspect}: the boolean operator joining
\verb|value is not None| and \verb|not value.isdigit()|.
\textit{What varies}: only the operator (\verb|or|~$\to$~\verb|and|);
the sub-conditions, the message, and the surrounding control flow are
invariant, isolating the operator as the cause.
\textit{Discerned}: the traceback chain
\verb|__init__| $\to$ setter $\to$ \verb|raise| exposes the property
mechanism; with \verb|"0701233455"| (digit, non-\verb|None|), the
\verb|or| is forced true regardless of \verb|isdigit()|.
\textit{Not yet discernible}: behaviour for \verb|value=None| (the
TODO's stated goal) and for non-digit strings --- only one of three
telling input values is tried.

\vspace{1em}

\noindent
\begin{minipage}[t]{0.48\linewidth}
\paragraph{Edit 2}
\begin{minted}{diff}
--- a/contacts.py
+++ b/contacts.py
@@ -32,7 +32,7 @@
         '''  sets phone number to a new value if it's a number'''
 
         # TODO: edit so value can be set to None
-        if  value is not None or not value.isdigit():
+        if value is not None and not value.isdigit():
             raise ValueError("Telefonnummret får bara innehålla siffror")
         self.__phone = value
\end{minted}
\end{minipage}
\hfill
\begin{minipage}[t]{0.48\linewidth}
\paragraph{Run 2}
Script: \texttt{contacts.py} \\
Result: Error \\
Exit code: 1

\begin{minted}{pytb}
KeyboardInterrupt
\end{minted}

\begin{minted}{text}
Välj ett alternativ (heltal):
1. Lägg till kontakt.
2. Sök efter kontakt.
3. Ändra en kontakt.
4. Skriv ut alla kontakter
5. Avsluta
  Namn: Daniel
  Tele: 0701233455
  Mail: dbosk@kth.se
  Adr : Någonannanstans
  Namn: Daniel
  Tele: 0701233455
  Mail: dbosk@kth.se
  Adr : Någonannanstans
  Namn: Kawsar
  Tele: 070
  Mail: kiaa@kth.se
  Adr : Någonstans
Traceback (most recent call last):
  File "/home/dbosk/teaching/prgi25/övn-2025-11-07/contacts.py", line 228, in <module>
    main()
    ~~~~^^
  File "/home/dbosk/teaching/prgi25/övn-2025-11-07/contacts.py", line 217, in main
    find_contact(contactbook)
    ~~~~~~~~~~~~^^^^^^^^^^^^^
  File "/home/dbosk/teaching/prgi25/övn-2025-11-07/contacts.py", line 135, in find_contact
    search_terms = input("Ange sökterm: ").strip()
                   ~~~~~^^^^^^^^^^^^^^^^^^
KeyboardInterrupt
\end{minted}
\end{minipage}

\paragraph*{Variation-theory note}
\textit{Reading the juxtaposition}: with the setter unblocked, execution
now reaches further and surfaces a \emph{different} defect --- a
\verb|while True| in \verb|find_contact| with no break condition,
escaped here only by \verb|KeyboardInterrupt|.
\textit{Discerned} (by contrast with Run 1): Edit 2's fix held; the
constructor no longer raises, and the new failure point lies elsewhere.
\textit{Not yet discernible}: whether the search results are correct.
The transcript captures stdout but not stdin, so the queries that
produced the listed matches are invisible.

\vspace{1em}

\noindent
\begin{minipage}[t]{0.48\linewidth}
\paragraph{Run 2}
Script: \texttt{contacts.py} \\
Result: Error \\
Exit code: 1

\begin{minted}{pytb}
KeyboardInterrupt
\end{minted}

\begin{minted}{text}
Välj ett alternativ (heltal):
1. Lägg till kontakt.
2. Sök efter kontakt.
3. Ändra en kontakt.
4. Skriv ut alla kontakter
5. Avsluta
  Namn: Daniel
  Tele: 0701233455
  Mail: dbosk@kth.se
  Adr : Någonannanstans
  Namn: Daniel
  Tele: 0701233455
  Mail: dbosk@kth.se
  Adr : Någonannanstans
  Namn: Kawsar
  Tele: 070
  Mail: kiaa@kth.se
  Adr : Någonstans
Traceback (most recent call last):
  File "/home/dbosk/teaching/prgi25/övn-2025-11-07/contacts.py", line 228, in <module>
    main()
    ~~~~^^
  File "/home/dbosk/teaching/prgi25/övn-2025-11-07/contacts.py", line 217, in main
    find_contact(contactbook)
    ~~~~~~~~~~~~^^^^^^^^^^^^^
  File "/home/dbosk/teaching/prgi25/övn-2025-11-07/contacts.py", line 135, in find_contact
    search_terms = input("Ange sökterm: ").strip()
                   ~~~~~^^^^^^^^^^^^^^^^^^
KeyboardInterrupt
\end{minted}
\end{minipage}
\hfill
\begin{minipage}[t]{0.48\linewidth}
\paragraph{Edit 3}
\begin{minted}{diff}
--- a/contacts.py
+++ b/contacts.py
@@ -132,7 +132,11 @@
     Hanterar menydelen för att söka efter kontakter.
     """
     while True:
-        search_terms = input("Ange sökterm: ").strip()
+        search_terms = input("Ange sökterm: ")
+        if not search_terms:
+            break
+        else:
+            search_terms = search_terms.strip()
 
         num_match_tuples = contactbook.find_contact(search_terms)
\end{minted}
\end{minipage}

\paragraph*{Variation-theory note}
\textit{Critical aspect}: how an interactive \verb|while True| returns
control to its caller.
\textit{What varies}: the input handling --- a single \verb|.strip()|
chain becomes a branch on emptiness with \verb|break|.
\textit{Discerned}: empty input is now a recognised exit signal; the
inner loop becomes escapable without \verb|Ctrl-C|.
\textit{Not yet discernible}: whether \verb|__lt__| orders ties
correctly (only two contacts exist), whether multi-term queries
accumulate match counts as intended, and whether
\verb|match_contact|'s exact-equality branch for phone numbers behaves
differently from the substring branches --- none of those dimensions
varies.

\vspace{1em}

\noindent
\begin{minipage}[t]{0.48\linewidth}
\paragraph{Edit 3}
\begin{minted}{diff}
--- a/contacts.py
+++ b/contacts.py
@@ -132,7 +132,11 @@
     Hanterar menydelen för att söka efter kontakter.
     """
     while True:
-        search_terms = input("Ange sökterm: ").strip()
+        search_terms = input("Ange sökterm: ")
+        if not search_terms:
+            break
+        else:
+            search_terms = search_terms.strip()
 
         num_match_tuples = contactbook.find_contact(search_terms)
\end{minted}
\end{minipage}
\hfill
\begin{minipage}[t]{0.48\linewidth}
\paragraph{Run 3}
Script: \texttt{contacts.py} \\
Result: Success \\
Exit code: 0

\begin{minted}{text}
Välj ett alternativ (heltal):
1. Lägg till kontakt.
2. Sök efter kontakt.
3. Ändra en kontakt.
4. Skriv ut alla kontakter
5. Avsluta
  Namn: Daniel
  Tele: 0701233455
  Mail: dbosk@kth.se
  Adr : Någonannanstans
Välj ett alternativ (heltal):
1. Lägg till kontakt.
2. Sök efter kontakt.
3. Ändra en kontakt.
4. Skriv ut alla kontakter
5. Avsluta
\end{minted}
\end{minipage}

\paragraph*{Variation-theory note}
\textit{Reading the juxtaposition}: the contrast Run~2~$\to$~Run~3
(with Edit 3 between) is exactly the contrast pattern variation
theory recommends --- same invocation, only the loop-exit logic
differs.
\textit{Discerned}: the menu now reaches \verb|Exit code: 0|; option 5
ends the program cleanly.
\textit{Not yet discernible}: the search returned only \emph{Daniel},
but the query is unrecorded --- a ``one expected match'' reading is
indistinguishable from a ``should have been two, returned one''
reading.

\vspace{1em}

\noindent
\begin{minipage}[t]{0.48\linewidth}
\paragraph{Run 3}
Script: \texttt{contacts.py} \\
Result: Success \\
Exit code: 0

\begin{minted}{text}
Välj ett alternativ (heltal):
1. Lägg till kontakt.
2. Sök efter kontakt.
3. Ändra en kontakt.
4. Skriv ut alla kontakter
5. Avsluta
  Namn: Daniel
  Tele: 0701233455
  Mail: dbosk@kth.se
  Adr : Någonannanstans
Välj ett alternativ (heltal):
1. Lägg till kontakt.
2. Sök efter kontakt.
3. Ändra en kontakt.
4. Skriv ut alla kontakter
5. Avsluta
\end{minted}
\end{minipage}
\hfill
\begin{minipage}[t]{0.48\linewidth}
\end{minipage}

\paragraph*{Variation-theory note}
This row has no contrasting partner: the right column is empty and the
left column repeats Run 3 unchanged. With nothing varying, no new
aspect can be discerned --- a useful reminder that
\emph{repetition is not variation}.

\vspace{1em}

\noindent
\begin{minipage}[t]{0.48\linewidth}
\end{minipage}
\hfill
\begin{minipage}[t]{0.48\linewidth}
\paragraph{Run 4}
Script: \texttt{contacts.py} \\
Result: Success \\
Exit code: 0

\begin{minted}{text}
Välj ett alternativ (heltal):
1. Lägg till kontakt.
2. Sök efter kontakt.
3. Ändra en kontakt.
4. Skriv ut alla kontakter
5. Avsluta
  Namn: Daniel
  Tele: 0701233455
  Mail: dbosk@kth.se
  Adr : Någonannanstans
Välj ett alternativ (heltal):
1. Lägg till kontakt.
2. Sök efter kontakt.
3. Ändra en kontakt.
4. Skriv ut alla kontakter
5. Avsluta
\end{minted}
\end{minipage}

\paragraph*{Variation-theory note}
Run 4 is identical to Run 3 in both inputs and outputs. Repeating a
successful run with the same inputs adds no new variation, hence no
new discernment. To learn more about the search and sort, one would
need to vary either the query (different terms, multi-term, exact
phone), the contact set (more rows, ties on different attributes), or
both.

\vspace{1em}

\noindent
\begin{minipage}[t]{0.48\linewidth}
\paragraph{Run 4}
Script: \texttt{contacts.py} \\
Result: Success \\
Exit code: 0

\begin{minted}{text}
Välj ett alternativ (heltal):
1. Lägg till kontakt.
2. Sök efter kontakt.
3. Ändra en kontakt.
4. Skriv ut alla kontakter
5. Avsluta
  Namn: Daniel
  Tele: 0701233455
  Mail: dbosk@kth.se
  Adr : Någonannanstans
Välj ett alternativ (heltal):
1. Lägg till kontakt.
2. Sök efter kontakt.
3. Ändra en kontakt.
4. Skriv ut alla kontakter
5. Avsluta
\end{minted}
\end{minipage}
\hfill
\begin{minipage}[t]{0.48\linewidth}
\paragraph{Edit 5}
\begin{minted}{diff}
--- a/contacts.py
+++ b/contacts.py
@@ -226,7 +226,9 @@
         elif choice == 5:
             return
         else:
-            raise NotImplementedError("Inga andra alternativ är implementerade.")
+            raise NotImplementedError(
+                "Inga andra alternativ är implementerade."
+            )
 
 if __name__ == "__main__":
     main()
\end{minted}
\end{minipage}

\paragraph*{Variation-theory note}
\textit{Reading the juxtaposition}: a deliberate \emph{non-variation}.
Edit 5 changes only whitespace inside the \verb|NotImplementedError|
call; no semantic dimension varies. If Run 5 differs from Run 4, the
cause must lie elsewhere (in the input choice). Formatting edits
paired with behavioural runs serve as invariance checks --- they
should leave behaviour unchanged.

\vspace{1em}

\noindent
\begin{minipage}[t]{0.48\linewidth}
\paragraph{Edit 5}
\begin{minted}{diff}
--- a/contacts.py
+++ b/contacts.py
@@ -226,7 +226,9 @@
         elif choice == 5:
             return
         else:
-            raise NotImplementedError("Inga andra alternativ är implementerade.")
+            raise NotImplementedError(
+                "Inga andra alternativ är implementerade."
+            )
 
 if __name__ == "__main__":
     main()
\end{minted}
\end{minipage}
\hfill
\begin{minipage}[t]{0.48\linewidth}
\paragraph{Run 5}
Script: \texttt{contacts.py} \\
Result: Error \\
Exit code: 1

\begin{minted}{pytb}
NotImplementedError: Inga andra alternativ är implementerade.
\end{minted}

\begin{minted}{text}
Välj ett alternativ (heltal):
1. Lägg till kontakt.
2. Sök efter kontakt.
3. Ändra en kontakt.
4. Skriv ut alla kontakter
5. Avsluta
Traceback (most recent call last):
  File "/home/dbosk/teaching/prgi25/övn-2025-11-07/contacts.py", line 234, in <module>
    main()
    ~~~~^^
  File "/home/dbosk/teaching/prgi25/övn-2025-11-07/contacts.py", line 229, in main
    raise NotImplementedError(
        "Inga andra alternativ är implementerade."
    )
NotImplementedError: Inga andra alternativ är implementerade.
\end{minted}
\end{minipage}

\paragraph*{Variation-theory note}
\textit{Critical aspect}: the \verb|else| arm of the menu dispatcher.
\textit{What varies} (vs.\ earlier runs): the menu choice --- an
unrecognised option. With Edit 5 invariant on semantics, the new
failure must come from the input dimension.
\textit{Discerned}: the catch-all correctly raises
\verb|NotImplementedError|.
\textit{Not yet discernible}: options 3 and 4 (\verb|pass| placeholders)
silently return to the menu --- behaviour distinct from option 5
(clean return) and from the catch-all (raise); none of the runs
separates these three outcomes. Likewise, non-numeric input (which
would raise \verb|ValueError| from \verb|int(...)|) is never tried.

\vspace{1em}


\section{How Dan-Claude did the variation theoretic analysis}

\inputminted{text}{opus-session.txt}

\end{document}
